% Original teaching diagram from the exact example in Chapter 19.
\begin{tikzpicture}[x=1mm,y=1mm,>=Stealth,font=\figurefont\fontsize{8}{10}\selectfont]
  \node[anchor=east] at (-4,20) {$\lambda$};
  \node[anchor=east] at (-4,0) {$1/\lambda$};
  \draw[->,FigureStroke] (0,20)--(82,20);
  \draw[->,FigureStroke] (0,0)--(82,0);
  % Both axes share the same linear scale, with one at x=40.
  \draw[FigureBlue,line width=1.15pt] (20,20)--(60,20);
  \draw[FigureGreen,line width=1.15pt] (24,0)--(66.667,0);
  \draw[FigureMuted,densely dashed,->] (20,19)--(66.667,1.1);
  \draw[FigureMuted,densely dashed,->] (60,19)--(24,1.1);
  \foreach \x/\t in {20/{3/4},40/1,60/{5/4}} {
    \draw[FigureStroke] (\x,19)--(\x,21);
    \node[above=1pt] at (\x,21) {$\t$};
  }
  \foreach \x/\t in {24/{4/5},40/1,66.667/{4/3}} {
    \draw[FigureStroke] (\x,-1)--(\x,1);
    \node[below=1pt] at (\x,-1) {$\t$};
  }
\end{tikzpicture}
